\documentclass[10pt,a4paper]{amsart}
\usepackage[margin=19mm]{geometry}
\usepackage{amsmath,amssymb,mathtools}
\usepackage[scr=rsfs,cal=euler]{mathalfa}
\usepackage{xfrac}
\usepackage[unicode,hidelinks,bookmarks=false]{hyperref}
\newcommand{\ie}{i.e.\ }
\newcommand{\cf}{cf.\ }
\newcommand{\resp}{respectively\ }
\newcommand{\Subsection}{\subsection}
\providecommand{\nobreakdash}{\nobreak\hbox{-}}
\setlength{\emergencystretch}{2em}
\multlinegap0pt
\usepackage{bm}
\usepackage{bbm}
\usepackage{dsfont}
\usepackage{xspace}
\usepackage{cancel}
\usepackage{mathtools}
\usepackage[shortlabels]{enumitem}
\usepackage{amsthm}
\makeatletter
\@ifundefined{assumption}{\newtheorem{assumption}{Assumption}}{}
\@ifundefined{thm}{\newtheorem{thm}{Theorem}}{}
\@ifundefined{cor}{\newtheorem{cor}[thm]{Corollary}}{}
\@ifundefined{lem}{\newtheorem{lem}[thm]{Lemma}}{}
\makeatother
\newcommand{\GDP}{G_D^\psi}
\newcommand{\LLL}{L^1(D,V(\de(x))dx)}
\newcommand{\calL}{\mathcal{L}^1}
\newcommand{\de}{\delta_D}
\newcommand{\diam}{\textrm{diam}}
\newcommand{\oa}[1]{\overline{a}_{#1}}
\newcommand{\od}[1]{\overline{\delta}_{#1}}
\newcommand{\ua}[1]{\underline{a}_{#1}}
\newcommand{\ud}[1]{\underline{\delta}_{#1}}
\title[Nonhomogeneous boundary condition for spectral non-local ope]{Nonhomogeneous boundary condition for spectral non-local operators}
\author{Ivan Bio\v{c}i\'c, Vanja Wagner}
\date{Source: arXiv:2601.21674v1 (CC BY 4.0)}
\begin{document}
\maketitle

\noindent\emph{Source and licence.} Nonhomogeneous boundary condition for spectral non-local operators. Version arXiv:2601.21674v1 (CC BY 4.0), distributed under the Creative Commons Attribution 4.0 International licence, as recorded in the arXiv licence metadata for this exact version and shown on the abstract page https://arxiv.org/abs/2601.21674v1. Full rights notice: licenses/arxiv-2601.21674-CC-BY-4.0.txt.\par\smallskip

\noindent\textit{Editorial scope.} Counted unit: the Thm whose source label is \texttt{t:conti-GDP-L1} in the article (source0.tex lines 644--647, source label \texttt{t:conti-GDP-L1}), delivered together with the article's own complete original proof of it (source0.tex lines 648--691, one \texttt{proof} environment of 44 lines). No mathematics is written, repaired or re-derived: statement and proof are the article's own text line for line. Required context retained uncounted: as:WSC (lines 236--242); U1 (lines 573--584); U4 (lines 573--584); eq:G(U) sharp bound (lines 586--591); c:bndry-G(U)-decay (lines 595--600); eq: U3' (lines 597--599); ap:l:rho (lines 1536--1539). Every definition, display and label of those lines is retained unchanged. Omitted from this package and not redistributed: 1--235, 243--572, 585--585, 592--594, 600--643, 692--1535, 1540--1720. Nothing inside a retained excerpt is omitted or altered: every retained body is delivered exactly as the article's lines are written, so no omission receipt of any kind accompanies this package. Citations: the retained excerpts carry no citation command at all, so no bibliography entry of the article is transcribed and no reference is redistributed. Reference adaptation: none was needed and none was made; every reference inside the delivered unit resolves inside it, so no unresolved reference or citation is produced. Printed slips kept literal: no printed word, symbol, hypothesis or number of the counted statement or of its proof was altered. Layout and class: the article's own class is \texttt{article} with packages including inputenc, fullpage, amsmath, amsthm, amsfonts, mathtools, xcolor, amssymb, enumerate, enumitem, comment, verbatim, hyperref, aligned-overset; the wrapper is the lane's pinned \texttt{amsart} editorial wrapper at 10pt on A4, which restates the theorem-like environments the retained text needs and adds the article's own macro definitions that the retained text needs (\texttt{\textbackslash GDP}, \texttt{\textbackslash LLL}, \texttt{\textbackslash calL}, \texttt{\textbackslash de}, \texttt{\textbackslash diam}, \texttt{\textbackslash oa}, \texttt{\textbackslash od}, \texttt{\textbackslash ua}, \texttt{\textbackslash ud}), with extra wrapper packages (bm, bbm, dsfont, xspace, cancel, mathtools, enumitem). The delivered numbering is the unit's own. Assets: no retained span contains a figure environment, a graphics-inclusion command, a PDF-page inclusion, a table or a TikZ picture; no image file is copied into this package, the package declares no asset dependency and no image file is redistributed with it. Licence: version arXiv:2601.21674v1 of this article is distributed under the Creative Commons Attribution 4.0 International licence, as recorded in the arXiv licence metadata for this exact version and shown on the abstract page https://arxiv.org/abs/2601.21674v1. A full rights notice is delivered at licenses/arxiv-2601.21674-CC-BY-4.0.txt. Characters: every retained excerpt is pure ASCII, so the delivered engine, XeLaTeX with its default Latin Modern text font, reports no missing character.
    \begin{assumption}{A}{1}\label{as:WSC}
        The exponent $\phi$ is a complete Bernstein functions, i.e.~the L\'evy measure $\mu$ has a completely monotone density denoted by $\mu(t)$, and $\phi$ satisfies the weak scaling condition at infinity: there exist $\ua{\phi},\oa{\phi}>0$ and $0< \ud{\phi}\le \od{\phi}<1$ such that
        \begin{align}\label{eq:WSC}
            \ua{\phi} \lambda^{\ud{\phi}} \phi(t)\le \phi(\lambda t)\le \oa{\phi}
            \lambda^{\od{\phi}}\phi(t),\quad \lambda\ge 1, \, t\ge 1.
        \end{align}
    \end{assumption}

		\begin{enumerate}[label=(\textbf{U\arabic*})]
			\item \label{U1}  
			$\int_{0}^1U(t)V(t)\, dt <\infty$;
			
			\item \label{U2} $U$ is almost non-increasing, i.e.~$\exists C>0$ such that $U(t)\le C U(s)$, $0<s\le t\le 1$;
			
			
			\item \label{U3} reverse doubling condition holds, i.e.~$\exists C>0$  
			such that $U(t)\le C U(2t)$, $t\in (0,1)$; 
			\item \label{U4} $U$ is bounded from
			above on $[c, \infty)$ for each $c >0$.
		\end{enumerate}

		\begin{equation}\label{eq:G(U) sharp bound}
    			\GDP\big(U(\de)\big)\asymp
    			\frac{1}{\de V(\de)\psi(V(\de)^{-2})}\int_0^{\de}U(t)V(t)\, dt+
    			V(\de)\left(1+\int_{\de}^{\diam D}
    			\frac{U(t)}{t V(t)\psi(V(t)^{-2})}\, dt\,\right).  
		\end{equation}

    \begin{cor}\label{c:bndry-G(U)-decay}
     The estimate \eqref{eq:G(U) sharp bound} holds for all functions $U:(0,\infty)\to(0,\infty)$ satisfying \ref{U1} and \ref{U4} such that for every $c>0$ there exists $C=C(c)>1$ such that 
        \begin{align}\label{eq: U3'}
            \frac{1}{C}\, U(t)\le U(s)\le C\, U(t),\quad   {c}^{-1} t \le s \le c\,t, \ s,t\in (0,1].  \tag{\textbf{U3'}}
        \end{align}
    \end{cor}

    \begin{lem}\label{ap:l:rho}
        The weight $\rho(\de)$ satisfies the conditions \ref{U1}, \eqref{eq: U3'} and \ref{U4}.
        
    \end{lem}

    \begin{thm}\label{t:conti-GDP-L1}
        It holds that $\GDP\big(\rho(\de)\big)\asymp V(\de)$, and the map $\GDP:\LLL\to \calL$
        is continuous.
    \end{thm}

    \begin{proof}
    First, assume that $\GDP\big(\rho(\de)\big)\asymp V(\de)$. Then the continuity of $\GDP$ follows by Fubini's theorem and the symmetry of $G_D^\psi(x,y)$, since
    \begin{align*}
        \|\GDP f\|_{\calL}&\le \int_D \GDP|f|(x)\rho(\de(x))dx=\int_D |f|(x)\GDP\big(\rho(\de)\big)(x)dx\\
        &\asymp \int_D |f|(x)V(\de(x))dx=\|f\|_{\LLL}.
    \end{align*}
 To prove $\GDP\big(\rho(\de)\big)\asymp V(\de)$ we treat each term in $\rho(\de)$ separately. By \ref{as:WSC} for $\phi$ and $\psi$ we have that \eqref{eq: U3'} holds for $t\mapsto V(t)^2\psi(V(t)^{-2})$, see Lemma \ref{ap:l:rho}. Hence, the formula \eqref{eq:G(U) sharp bound} implies
    \begin{align*}
        \GDP\big(V(\de)^2\psi(V(\de)^{-2})\big)&\asymp \frac{1}{\de V(\de)\psi(V(\de)^{-2})}\int\limits_0^{\de}V(t)^3\psi(V(t)^{-2})\, dt\\
    			&\hspace{-1em}+
    			V(\de)+V(\de)\int\limits_{\de}^{\diam D}
    			\frac{V(t)}{t}\, dt\,.  
    \end{align*}
    Since $r\mapsto r^2\psi(r^{-2})$ is non-decreasing, and $V$ satisfies the weak scaling condition, the first term is bounded from above by $\frac{V(\de)}{\de}\int_0^{\de} V(t)dt\asymp  V(\de)^2\lesssim V(\de).$
     
    For the third term, we can again use the weak scaling property of $V$ to get $V(\de)\int_{\de}^{\diam D}
    			\frac{V(t)}{t}\, dt\lesssim  V(\de)\int_{0}^{\diam D}
    			\frac{V(t)}{t}\, dt\lesssim V(\de).$
    Therefore, $\GDP\big(V(\de)^2\psi(V(\de)^{-2})\big)\asymp V(\de)$.

    The second term of $\rho(\de)$ also satisfies \eqref{eq: U3'}, see Lemma \ref{ap:l:rho}.  Thus, we can apply Corollary \ref{c:bndry-G(U)-decay} to get
    \begin{align*}
        \GDP\left(V(\de)\int_{V(\de)}^1\psi(s^{-2})ds\right)&\asymp \frac{1}{\de(x)V(\de)\psi(V(\de)^{-2})}\int_0^{\de(x)}V(t)^2\int_{V(t)}^1\psi(s^{-2})ds\, dt\\
    			&\hspace{-1em}+
    			V(\de)+V(\de)\int_{\de}^{\diam D}
    			\frac{\int_{V(t)}^1\psi(s^{-2})ds}{t\psi(V(t)^{-2})}\, dt\,.  
    \end{align*}
    Note that by the weak scaling of $\psi$ we have
    \begin{align}\label{eq:1802}
        \int_{V(t)}^1\psi(s^{-2})ds\lesssim\begin{cases}
            V(t)^{2\ud{\psi}} \psi(V(t)^{-2}),& \ud{\psi}<\frac12,\\
            V(t)\psi(V(t)^{-2})\log\big(1/t\big),& \ud{\psi}=\frac12,\\
            V(t)\psi(V(t)^{-2}),& \ud{\psi}\ge \frac12.
        \end{cases} \lesssim V(t)^{\eta}\psi(V(t)^{-2})
    \end{align}
    for some $\eta\in (0,1)$. Hence, it follows that
    \begin{align*}
       \GDP\left(V(\de)\int_{V(\de)}^1\psi(s^{-2})ds\right)&\lesssim V(\de)^{1+\eta}+
    			V(\de)\left(1+\int_{\de}^{\diam D}
    			\frac{V(t)^{\eta}}{t}\, dt\,\right)\lesssim V(\de),
    \end{align*}
    where for the first term we used that $r\mapsto r^2\psi(^{-2})$ is non-decreasing, and that $V$ is increasing, and for the third we note that $\int_{0}^{\diam D}\frac{V(t)^{\eta}}{t}\, dt<\infty$ since $V(t)\lesssim t^{\ud{\phi}}$ by the weak scaling at zero and $\eta \ud{\phi}<1$.
    This finishes the proof of $\GDP\rho(\de)\asymp V(\de)$.
    \end{proof}

\end{document}
